\subsection{General Linear Elliptic Problems}\label{sec:general-linear-elliptic-problems}

So far in this chapter, we have discussed only Dirichlet boundary conditions for elliptic problems. In this section we shall discuss a few of the other boundary conditions that arise in physical and mathematical problems. Since physical problems present us with a wide variety of boundary conditions for consideration our discussion will not be exhaustive.

\subsubsection{The Neumann Problem}\label{subsec:neumann-problem}

After the Dirichlet problem, the second most common and important elliptic boundary-value problem is the Neumann problem.

\begin{Definition}
\label{def:classical-neumann-solution}
Let $\Omega\subset\mathbb{R}^n$ be a bounded domain with $C^1$ boundary and suppose $f\in C_b(\Omega)$ is given. A function\[u\in C_b^2(\Omega)\cap C_b^1(\overline\Omega)\]is a \textbf{classical solution} of the Neumann problem if
\begin{equation}\label{eq:neumann-classical-divergence-equation}
L(\bm{x},D)u:=\sum_{i,j=1}^n\frac{\partial}{\partial x_i}a_{ij}(\bm{x})\frac{\partial u}{\partial x_j}+\sum_{i=1}^nb_i(\bm{x})\frac{\partial u}{\partial x_i}+c(\bm{x})u=f
\end{equation}
in $\Omega$; and
\begin{equation}\label{eq:neumann-homogeneous-conormal-boundary-condition}
\sum_{i,j=1}^na_{ij}(\bm{x})\frac{\partial u}{\partial x_j}\eta_i(\bm{x})=0
\end{equation}
on $\partial\Omega$ where $\bm{\eta}(\bm{x})$ is the unit outward normal to $\partial\Omega$ at $\bm{x}$.
\end{Definition}

As with the Dirichlet problem, we can define a strong solution of the Neumann problem.

\begin{Definition}
\label{def:strong-neumann-solution}
Let $\Omega\subset\mathbb{R}^n$ be a bounded domain with a $C^1$ boundary and suppose $f\in L^2(\Omega)$ is given. A function\[u\in H^2(\Omega)\]is a \textbf{strong solution} of the Neumann problem if \eqref{eq:neumann-classical-divergence-equation} holds in $L^2(\Omega)$ and \eqref{eq:neumann-homogeneous-conormal-boundary-condition} holds in the sense of trace on $\partial\Omega$.
\end{Definition}

In order to state the Neumann problem in weak form we proceed as before and use integration by parts to create a bilinear form from the differential operator $L$ and the boundary conditions. Note that for any $\phi$ and $u$ in $H^2(\Omega)$ we have
\begin{equation}\label{eq:neumann-integration-by-parts-boundary-form}
\begin{aligned}
\int_\Omega\phi Lu\,\mathrm{d}\bm{x}&=B[\phi,u]\\
&\quad+\sum_{i,j=1}^n\int_{\partial\Omega}a_{ij}(\bm{x})\frac{\partial u}{\partial x_j}(\bm{x})\eta_i(\bm{x})\phi(\bm{x})\,\mathrm{d}\bm{x},
\end{aligned}
\end{equation}
where the bilinear form $B$ is defined in \eqref{eq:second-order-garding-bilinear-form}. Thus, if $u$ satisfies the boundary condition \eqref{eq:neumann-homogeneous-conormal-boundary-condition}, then we have
\begin{equation}\label{eq:neumann-natural-variational-integration-identity}
\int_\Omega\phi Lu\,\mathrm{d}\bm{x}=B[\phi,u]
\end{equation}
for every $\phi\in H^1(\Omega)$. In fact, we take this as the definition of a weak solution of the Neumann problem.

\begin{Definition}
\label{def:weak-neumann-solution}
Let $\Omega\subset\mathbb{R}^n$ be a bounded domain and suppose $f\in L^2(\Omega)$ is given. A function\[u\in H^1(\Omega)\]is a \textbf{weak solution} of the Neumann problem if
\begin{equation}\label{eq:weak-neumann-variational-equation}
B[v,u]=(f,v)
\end{equation}
for every $v\in H^1(\Omega)$.
\end{Definition}

A few comments are in order.\begin{enumerate}\item Of course, we have constructed things so that every strong solution of the Neumann problem is also a weak solution.\item In the construction of the weak form, the boundary conditions ``disappear\textquotesingle\textquotesingle; i.e., condition \eqref{eq:neumann-homogeneous-conormal-boundary-condition} does not appear explicitly in either the bilinear form or the space of admissible functions. Because of this, Neumann conditions are referred to as \emph{natural boundary conditions}. In fact, since the trace theorem does not guarantee the existence of normal derivatives of $H^1(\Omega)$ functions, it does not necessarily make sense to evaluate the boundary condition \eqref{eq:neumann-homogeneous-conormal-boundary-condition} on a weak solution of \eqref{eq:weak-neumann-variational-equation}.\item We have assumed that the data $f$ is in the space $L^2(\Omega)$. This can be weakened, for instance by taking $f\in L^2(S)$ where $S$ is a smooth surface contained in $\Omega$. However, we cannot take arbitrary data in $H^{-1}(\Omega)$ as we did for the Dirichlet problem.\end{enumerate}
